\documentclass[11pt]{article}
\usepackage{mathblatt}

\begin{document}
\blattkopf{Prozentrechnung}{Fokus Prozentsatz · Lösungen}
\begleitteil

\einheitenkopf*{Blatt 0 · Das kannst du schon}

Blatt 0: Nr. 1 Grundvorstellung Anteil · Nr. 2 bis 5 die Fertigkeiten für Nr. 6 bis 16.

\erg{1}{a) die Hälfte \quad b) drei Viertel \quad c) zwei von vier Kästchen gefärbt \quad d) ein von zehn Kästchen gefärbt}
\erg{2}{a) $\frac{15}{100}$ \quad b) $\frac{28}{100}$ \quad c) $\frac{12}{40} = \frac{3}{10} = \frac{30}{100}$}
\erg{3}{a) 45 \quad b) 6 \quad c) $0{,}07$}
\erg{4}{a) 1 Becher $0{,}60\,€$; 100 Becher $60{,}00\,€$ \quad b) 1 Sticker $0{,}12\,€$; 100 Sticker $12{,}00\,€$}
\erg{5}{a) $4{,}2$ \quad b) $17{,}6$ \quad c) $7{,}0$}

\einheitenkopf*{Fokus · Prozentsatz berechnen}

\erg{6}{a) 24 Kinder \quad b) 50 Plätze \quad c) 15 Würfe \quad d) 200 Zuschauer \quad e) 18 und 30 zusammen (48 Kekse)}
\erg{7}{a) neun Striche, Abstand ein Zehntel der Länge \quad b) drei Striche, Abstand ein Viertel der Länge \quad c) Pfeil bei $50\,\%$, also in der Mitte}
\erg{8}{a) $40\,\%$ \quad b) $70\,\%$ \quad c) $90\,\%$ \quad d) $20\,\%$}
\erg{9}{a) $65\,\%$ \quad b) $15\,\%$ \quad c) $85\,\%$}
\erg{10}{a) $23\,\%$ \quad b) $8\,\%$ \quad c) $60\,\%$ \quad d) $\frac{3}{10} = \frac{30}{100} = 30\,\%$ \quad e) $\frac{9}{20} = \frac{45}{100} = 45\,\%$ \quad f) $\frac{19}{25} = \frac{76}{100} = 76\,\%$ \quad g) $\frac{21}{50} = \frac{42}{100} = 42\,\%$}
\erg{11}{a) $12 : 48 = 0{,}25 = 25\,\%$ \quad b) $27 : 60 = 0{,}45 = 45\,\%$ \quad c) $33 : 55 = 0{,}6 = 60\,\%$ \quad d) $5 : 12 = 0{,}4166\ldots \approx 41{,}7\,\%$ \quad e) $17 : 90 = 0{,}1888\ldots \approx 18{,}9\,\%$}
\erg{12}{mehrere Antworten richtig, z.\,B. a) 3 von 6 \quad b) 5 von 20 \quad c) 2 von 5. Probe: Teil geteilt durch Ganzes ergibt den Prozentsatz.}
\erg{13}{a) Ganzes 25; $13 : 25 = 0{,}52 = 52\,\%$ \quad b) Ganzes 40; $26 : 40 = 0{,}65 = 65\,\%$ \quad c) Ganzes $7 + 19 = 26$; $7 : 26 = 0{,}2692\ldots \approx 26{,}9\,\%$}
\erg{14}{a) Jonas hat Teil und Ganzes vertauscht; er teilt das Ganze durch den Teil. Richtig: $12 : 30 = 0{,}4 = 40\,\%$ \quad b) $16 : 64 = 0{,}25 = 25\,\%$}
\erg{15}{a) Der Prozentsatz sagt, wie viele Hundertstel des Ganzen der Teil ist; darum ist das Ganze der Divisor. \quad b) Gleich groß: 8 von 20 und 12 von 30 lassen sich beide zu 2 von 5 kürzen.}
\erg{16}{a) $1{,}5 : 6{,}4 = 0{,}2343\ldots \approx 23{,}4\,\%$ \quad b) nicht im Tor: $60 - 48 = 12$; $12 : 60 = 0{,}2 = 20\,\%$}

\end{document}
